% !TEX root = thesis.tex
\chapter{Introduction}
\label{chap:introduction}

Large generative models have significantly advanced the capabilities of Artificial
Intelligence (AI) systems in areas such as natural-language processing, code generation,
question answering, and reasoning. Much of this progress is based on the Transformer
architecture, whose attention mechanism enables models to construct context-aware
representations of input sequences \cite{vaswani2017attention}. Increasing the number of
model parameters and the amount of training data has further improved the range of tasks
that these models can perform, as demonstrated by model families such as GPT and LLaMA
\cite{brown2020gpt3,touvron2023llama}.

However, the capabilities of these models are accompanied by considerable inference
requirements. Depending on the model size, sequence length, numerical precision, and
deployment hardware, generating an output typically requires substantial computation,
memory capacity, memory bandwidth, and energy. These requirements are especially
restrictive on edge devices such as smartphones, embedded systems, and other platforms
that operate with limited memory, processing capacity, and power budgets. Consequently,
deploying large generative models outside specialised data centres often requires
optimisations at the model, algorithmic, and system levels
\cite{bai2024beyond,zhou2024efficientinf}.

A clear understanding of the inference process is necessary before these optimisation
methods can be evaluated. In decoder-based language models, inference consists of an
initial prefill phase followed by sequential autoregressive decoding. The two phases exhibit
different resource characteristics: prefill processes the prompt largely in parallel and can
be compute-intensive, whereas decoding repeatedly accesses the model parameters and the
growing \ac{kv} cache for every generated token. The dominant bottleneck may
therefore vary between computation, memory capacity, memory bandwidth, latency, and
energy depending on the workload and target platform
\cite{yuan2024llminference,zhou2024efficientinf}.

Against this background, this report addresses the following research questions:

\begin{itemize}
    \item Which operations and data structures constitute the main resource bottlenecks
    during Transformer inference?

    \item How do model compression, inference-level optimisation, decoding acceleration,
    and system-level runtimes address these bottlenecks?

    \item What trade-offs determine whether a particular optimisation is suitable for
    resource-constrained edge devices or server-based deployment?
\end{itemize}

To answer these questions, the report first introduces the relevant components of the
Transformer architecture and examines the prefill and decode phases of inference. It then
surveys model-compression techniques, including quantisation, pruning, and knowledge
distillation; inference-level methods for attention and KV-cache optimisation; decoding
and scheduling techniques such as speculative decoding and continuous batching; and
system-level runtimes including \texttt{llama.cpp}, \ac{mlc-llm}, and vLLM. The methods are
considered according to the bottlenecks they address, their resource--accuracy trade-offs,
and their suitability for different deployment environments.







